% Seminar 3 - Page 7. Fragment for
% % Seminar 3 - Page 7. Fragment for
% % Seminar 3 - Page 7. Fragment for
% % Seminar 3 - Page 7. Fragment for
% \input{seminars/pages/seminar3-pagina7-en}; no preamble or page break.
% EDITORIAL NOTES:
% - This page continues the examples from pages 5-6. The visible labels
%   (ii), (iii), (iv) are retained even where neighbouring pages differ.
% - The two triangular diagrams have been redrawn using tikz-cd.
% - For the argument map on C*, the notes suppress the choice of argument;
%   normalisation z/|z| is used to express the indicated correspondence.
% - In (iv), the handwritten reflexivity line says "0 in R" rather than
%   the needed "0 in Q". It is retained in the displayed source argument;
%   the discrepancy is recorded here, not silently corrected.
% - The sketch involving a map R -> R/~ and f(x)=f(y) iff x-y in Q is
%   visibly struck through. It is noted here but not asserted as a map.
% - R/~ ~ R is a set-theoretic assertion, not a canonical identification.
% - A handwritten sentence after p:R -> R/~ appears to say that R itself
%   is not a system of representatives; the reading is somewhat unclear.

\subsection*{Seminar 3 - Page 7}
\begingroup
\setlength{\abovedisplayskip}{5pt}
\setlength{\belowdisplayskip}{5pt}
\setlength{\abovedisplayshortskip}{3pt}
\setlength{\belowdisplayshortskip}{3pt}

\noindent\textbf{(ii) Equal-modulus classes (continued).}
The canonical projection $p$ and the modulus map $f$ give the diagram
\[
\begin{tikzcd}[column sep=large,row sep=small]
 \mathbb C\arrow[r,"p"]\arrow[d,"f"'] & \mathbb C/{\sim}\\
 {[0,+\infty)}\arrow[ur,dashed,"\overline f"'] &
\end{tikzcd}
\qquad f(z)=|z|.
\]
Here $\sim=R_f$, that is,
\[
 x\sim y\iff f(x)=f(y).
\]

\noindent\textbf{(iii) Equal-argument classes (continued).}
For $\mathbb C^*=\mathbb C\setminus\{0\}$ and the unit circle $S'$:
\[
\begin{tikzcd}[column sep=large,row sep=small]
 \mathbb C^*\arrow[r,"p"]\arrow[d,"f"'] & \mathbb C^*/{\sim}\\
 S'\arrow[ur,dashed,"\overline f"'] &
\end{tikzcd}
\]
Writing $z=r(\cos t+i\sin t)$, the notes give
\[
 f(z)=f\bigl(r(\cos t+i\sin t)\bigr)=\cos t+i\sin t.
\]

\noindent\textbf{(iv) Equivalence modulo $\mathbb Q$ on $\mathbb R$.}
Define
\[
 x\sim y\overset{\mathrm{def}}\iff x-y\in\mathbb Q,
 \qquad \mathbb R/{\sim}\cong\mathbb R.
\]
Check the three properties:
\[
\begin{aligned}
\text{(R)}\quad &x\sim x\quad(0\in\mathbb R),\\
\text{(S)}\quad &x-y\in\mathbb Q\Longrightarrow y-x\in\mathbb Q,\\
\text{(T)}\quad &x-y\in\mathbb Q,\ y-z\in\mathbb Q
      \Longrightarrow x-z\in\mathbb Q.
\end{aligned}
\]
The projection is $p:\mathbb R\to\mathbb R/{\sim}$.
The notes remark that $\mathbb R$ itself is not a system of representatives.
For a system $S=(x_i)_{i\in I}$ of representatives,
\[
 p|_S:S\longrightarrow\mathbb R/{\sim}\quad\text{is bijective}.
\]
Since
\[
 \mathbb R=\bigcup_{i\in I}\widehat{x_i},\qquad
 \widehat{x_i}=\{y\in\mathbb R\mid y-x_i\in\mathbb Q\}
   =x_i+\mathbb Q\cong\mathbb Q,
\]
each $\widehat{x_i}$ is countable. If $I$ were countable,
$\mathbb R$ would be countable, a contradiction. Hence $I$ is uncountable.

\endgroup
; no preamble or page break.
% EDITORIAL NOTES:
% - This page continues the examples from pages 5-6. The visible labels
%   (ii), (iii), (iv) are retained even where neighbouring pages differ.
% - The two triangular diagrams have been redrawn using tikz-cd.
% - For the argument map on C*, the notes suppress the choice of argument;
%   normalisation z/|z| is used to express the indicated correspondence.
% - In (iv), the handwritten reflexivity line says "0 in R" rather than
%   the needed "0 in Q". It is retained in the displayed source argument;
%   the discrepancy is recorded here, not silently corrected.
% - The sketch involving a map R -> R/~ and f(x)=f(y) iff x-y in Q is
%   visibly struck through. It is noted here but not asserted as a map.
% - R/~ ~ R is a set-theoretic assertion, not a canonical identification.
% - A handwritten sentence after p:R -> R/~ appears to say that R itself
%   is not a system of representatives; the reading is somewhat unclear.

\subsection*{Seminar 3 - Page 7}
\begingroup
\setlength{\abovedisplayskip}{5pt}
\setlength{\belowdisplayskip}{5pt}
\setlength{\abovedisplayshortskip}{3pt}
\setlength{\belowdisplayshortskip}{3pt}

\noindent\textbf{(ii) Equal-modulus classes (continued).}
The canonical projection $p$ and the modulus map $f$ give the diagram
\[
\begin{tikzcd}[column sep=large,row sep=small]
 \mathbb C\arrow[r,"p"]\arrow[d,"f"'] & \mathbb C/{\sim}\\
 {[0,+\infty)}\arrow[ur,dashed,"\overline f"'] &
\end{tikzcd}
\qquad f(z)=|z|.
\]
Here $\sim=R_f$, that is,
\[
 x\sim y\iff f(x)=f(y).
\]

\noindent\textbf{(iii) Equal-argument classes (continued).}
For $\mathbb C^*=\mathbb C\setminus\{0\}$ and the unit circle $S'$:
\[
\begin{tikzcd}[column sep=large,row sep=small]
 \mathbb C^*\arrow[r,"p"]\arrow[d,"f"'] & \mathbb C^*/{\sim}\\
 S'\arrow[ur,dashed,"\overline f"'] &
\end{tikzcd}
\]
Writing $z=r(\cos t+i\sin t)$, the notes give
\[
 f(z)=f\bigl(r(\cos t+i\sin t)\bigr)=\cos t+i\sin t.
\]

\noindent\textbf{(iv) Equivalence modulo $\mathbb Q$ on $\mathbb R$.}
Define
\[
 x\sim y\overset{\mathrm{def}}\iff x-y\in\mathbb Q,
 \qquad \mathbb R/{\sim}\cong\mathbb R.
\]
Check the three properties:
\[
\begin{aligned}
\text{(R)}\quad &x\sim x\quad(0\in\mathbb R),\\
\text{(S)}\quad &x-y\in\mathbb Q\Longrightarrow y-x\in\mathbb Q,\\
\text{(T)}\quad &x-y\in\mathbb Q,\ y-z\in\mathbb Q
      \Longrightarrow x-z\in\mathbb Q.
\end{aligned}
\]
The projection is $p:\mathbb R\to\mathbb R/{\sim}$.
The notes remark that $\mathbb R$ itself is not a system of representatives.
For a system $S=(x_i)_{i\in I}$ of representatives,
\[
 p|_S:S\longrightarrow\mathbb R/{\sim}\quad\text{is bijective}.
\]
Since
\[
 \mathbb R=\bigcup_{i\in I}\widehat{x_i},\qquad
 \widehat{x_i}=\{y\in\mathbb R\mid y-x_i\in\mathbb Q\}
   =x_i+\mathbb Q\cong\mathbb Q,
\]
each $\widehat{x_i}$ is countable. If $I$ were countable,
$\mathbb R$ would be countable, a contradiction. Hence $I$ is uncountable.

\endgroup
; no preamble or page break.
% EDITORIAL NOTES:
% - This page continues the examples from pages 5-6. The visible labels
%   (ii), (iii), (iv) are retained even where neighbouring pages differ.
% - The two triangular diagrams have been redrawn using tikz-cd.
% - For the argument map on C*, the notes suppress the choice of argument;
%   normalisation z/|z| is used to express the indicated correspondence.
% - In (iv), the handwritten reflexivity line says "0 in R" rather than
%   the needed "0 in Q". It is retained in the displayed source argument;
%   the discrepancy is recorded here, not silently corrected.
% - The sketch involving a map R -> R/~ and f(x)=f(y) iff x-y in Q is
%   visibly struck through. It is noted here but not asserted as a map.
% - R/~ ~ R is a set-theoretic assertion, not a canonical identification.
% - A handwritten sentence after p:R -> R/~ appears to say that R itself
%   is not a system of representatives; the reading is somewhat unclear.

\subsection*{Seminar 3 - Page 7}
\begingroup
\setlength{\abovedisplayskip}{5pt}
\setlength{\belowdisplayskip}{5pt}
\setlength{\abovedisplayshortskip}{3pt}
\setlength{\belowdisplayshortskip}{3pt}

\noindent\textbf{(ii) Equal-modulus classes (continued).}
The canonical projection $p$ and the modulus map $f$ give the diagram
\[
\begin{tikzcd}[column sep=large,row sep=small]
 \mathbb C\arrow[r,"p"]\arrow[d,"f"'] & \mathbb C/{\sim}\\
 {[0,+\infty)}\arrow[ur,dashed,"\overline f"'] &
\end{tikzcd}
\qquad f(z)=|z|.
\]
Here $\sim=R_f$, that is,
\[
 x\sim y\iff f(x)=f(y).
\]

\noindent\textbf{(iii) Equal-argument classes (continued).}
For $\mathbb C^*=\mathbb C\setminus\{0\}$ and the unit circle $S'$:
\[
\begin{tikzcd}[column sep=large,row sep=small]
 \mathbb C^*\arrow[r,"p"]\arrow[d,"f"'] & \mathbb C^*/{\sim}\\
 S'\arrow[ur,dashed,"\overline f"'] &
\end{tikzcd}
\]
Writing $z=r(\cos t+i\sin t)$, the notes give
\[
 f(z)=f\bigl(r(\cos t+i\sin t)\bigr)=\cos t+i\sin t.
\]

\noindent\textbf{(iv) Equivalence modulo $\mathbb Q$ on $\mathbb R$.}
Define
\[
 x\sim y\overset{\mathrm{def}}\iff x-y\in\mathbb Q,
 \qquad \mathbb R/{\sim}\cong\mathbb R.
\]
Check the three properties:
\[
\begin{aligned}
\text{(R)}\quad &x\sim x\quad(0\in\mathbb R),\\
\text{(S)}\quad &x-y\in\mathbb Q\Longrightarrow y-x\in\mathbb Q,\\
\text{(T)}\quad &x-y\in\mathbb Q,\ y-z\in\mathbb Q
      \Longrightarrow x-z\in\mathbb Q.
\end{aligned}
\]
The projection is $p:\mathbb R\to\mathbb R/{\sim}$.
The notes remark that $\mathbb R$ itself is not a system of representatives.
For a system $S=(x_i)_{i\in I}$ of representatives,
\[
 p|_S:S\longrightarrow\mathbb R/{\sim}\quad\text{is bijective}.
\]
Since
\[
 \mathbb R=\bigcup_{i\in I}\widehat{x_i},\qquad
 \widehat{x_i}=\{y\in\mathbb R\mid y-x_i\in\mathbb Q\}
   =x_i+\mathbb Q\cong\mathbb Q,
\]
each $\widehat{x_i}$ is countable. If $I$ were countable,
$\mathbb R$ would be countable, a contradiction. Hence $I$ is uncountable.

\endgroup
; no preamble or page break.
% EDITORIAL NOTES:
% - This page continues the examples from pages 5-6. The visible labels
%   (ii), (iii), (iv) are retained even where neighbouring pages differ.
% - The two triangular diagrams have been redrawn using tikz-cd.
% - For the argument map on C*, the notes suppress the choice of argument;
%   normalisation z/|z| is used to express the indicated correspondence.
% - In (iv), the handwritten reflexivity line says "0 in R" rather than
%   the needed "0 in Q". It is retained in the displayed source argument;
%   the discrepancy is recorded here, not silently corrected.
% - The sketch involving a map R -> R/~ and f(x)=f(y) iff x-y in Q is
%   visibly struck through. It is noted here but not asserted as a map.
% - R/~ ~ R is a set-theoretic assertion, not a canonical identification.
% - A handwritten sentence after p:R -> R/~ appears to say that R itself
%   is not a system of representatives; the reading is somewhat unclear.

\subsection*{Seminar 3 - Page 7}
\begingroup
\setlength{\abovedisplayskip}{5pt}
\setlength{\belowdisplayskip}{5pt}
\setlength{\abovedisplayshortskip}{3pt}
\setlength{\belowdisplayshortskip}{3pt}

\noindent\textbf{(ii) Equal-modulus classes (continued).}
The canonical projection $p$ and the modulus map $f$ give the diagram
\[
\begin{tikzcd}[column sep=large,row sep=small]
 \mathbb C\arrow[r,"p"]\arrow[d,"f"'] & \mathbb C/{\sim}\\
 {[0,+\infty)}\arrow[ur,dashed,"\overline f"'] &
\end{tikzcd}
\qquad f(z)=|z|.
\]
Here $\sim=R_f$, that is,
\[
 x\sim y\iff f(x)=f(y).
\]

\noindent\textbf{(iii) Equal-argument classes (continued).}
For $\mathbb C^*=\mathbb C\setminus\{0\}$ and the unit circle $S'$:
\[
\begin{tikzcd}[column sep=large,row sep=small]
 \mathbb C^*\arrow[r,"p"]\arrow[d,"f"'] & \mathbb C^*/{\sim}\\
 S'\arrow[ur,dashed,"\overline f"'] &
\end{tikzcd}
\]
Writing $z=r(\cos t+i\sin t)$, the notes give
\[
 f(z)=f\bigl(r(\cos t+i\sin t)\bigr)=\cos t+i\sin t.
\]

\noindent\textbf{(iv) Equivalence modulo $\mathbb Q$ on $\mathbb R$.}
Define
\[
 x\sim y\overset{\mathrm{def}}\iff x-y\in\mathbb Q,
 \qquad \mathbb R/{\sim}\cong\mathbb R.
\]
Check the three properties:
\[
\begin{aligned}
\text{(R)}\quad &x\sim x\quad(0\in\mathbb R),\\
\text{(S)}\quad &x-y\in\mathbb Q\Longrightarrow y-x\in\mathbb Q,\\
\text{(T)}\quad &x-y\in\mathbb Q,\ y-z\in\mathbb Q
      \Longrightarrow x-z\in\mathbb Q.
\end{aligned}
\]
The projection is $p:\mathbb R\to\mathbb R/{\sim}$.
The notes remark that $\mathbb R$ itself is not a system of representatives.
For a system $S=(x_i)_{i\in I}$ of representatives,
\[
 p|_S:S\longrightarrow\mathbb R/{\sim}\quad\text{is bijective}.
\]
Since
\[
 \mathbb R=\bigcup_{i\in I}\widehat{x_i},\qquad
 \widehat{x_i}=\{y\in\mathbb R\mid y-x_i\in\mathbb Q\}
   =x_i+\mathbb Q\cong\mathbb Q,
\]
each $\widehat{x_i}$ is countable. If $I$ were countable,
$\mathbb R$ would be countable, a contradiction. Hence $I$ is uncountable.

\endgroup
